% ============================================================
%  EXAMEN TEMA 2: NÚMEROS ENTEROS Y SUS OPERACIONES
%  Curso Introductorio de Matemáticas — 1er Año
%  Autor: Ing. Gabriel Astudillo
%  Sesiones cubiertas: Semana 2 (clases 5 a 8)
%  Nota: enunciado limpio (sin pistas ni desarrollo). Las
%  soluciones están en la "Clave de Respuestas" al final.
% ============================================================

\documentclass[12pt, letterpaper]{article}

% ── Paquetes ─────────────────────────────────────────────────

\usepackage{mathwizards-guia}
\mwguia{Examen — Tema 2: Números Enteros}{Ing. Gabriel Astudillo $\cdot$ Curso 1er Año}

% ── Comandos propios de este examen ───────────────────
\newcommand{\Z}{\mathbb{Z}}

\begin{document}
% ============================================================

% ── Portada / Título ─────────────────────────────────────────
\mwportada{Examen del Tema 2}{Los Números Enteros y sus Operaciones}{Suma, Resta, Multiplicación y División en $\Z$ con Jerarquía}

\vspace{4pt}
\begin{cuadroinfo}
  \small
  \textbf{Nombre:} \rule{0.55\linewidth}{0.4pt} \quad \textbf{Fecha:} \rule{0.15\linewidth}{0.4pt} \\[6pt]
  \textbf{Tiempo sugerido:} 60 minutos \quad$\bullet$\quad \textbf{Puntaje total:} 100 puntos \\[4pt]
  \textbf{Instrucciones:}
  \begin{enumerate}[leftmargin=*]
    \item Lee cada ejercicio con calma antes de resolverlo.
    \item Muestra \emph{todos} tus pasos: se evalúa cómo llegas a la respuesta, no solo la respuesta.
    \item Aplica siempre la \emph{regla de los signos} y la \emph{jerarquía de las operaciones}.
    \item En las operaciones combinadas, resuelve de adentro hacia afuera (paréntesis, corchetes, llaves).
    \item No se permite el uso de calculadora.
  \end{enumerate}
\end{cuadroinfo}

\vspace{8pt}

% ============================================================
\section{Suma y Resta en $\Z$}
% ============================================================

\begin{enumerate}[leftmargin=*]
  \item Calcule el valor exacto de:
        \[
          |-12 + 5| - |-8 - 3| + |-(-7)|
        \]

  \item Resuelva, simplificando primero los signos dobles:
        \[
          -(3 - 9) - [-(5 - 12) + 4]
        \]

  \item Obtenga el resultado de la siguiente cadena de signos:
        \[
          -(-15) + (-8) - (+6) - (-(-4))
        \]

  \item Encuentre el número entero que falta para que la igualdad sea verdadera:
        \[
          -10 + \_\_\_ - (-5) = -18
        \]
\end{enumerate}

\newpage

% ============================================================
\section{Multiplicación y División en $\Z$}
% ============================================================

\begin{enumerate}[leftmargin=*]
  \item Calcule:
        \[
          (-4) \times (-6) \div (-3) \times 2
        \]

  \item Calcule:
        \[
          (-9) \times 8 \div (-6) \div (-4)
        \]

  \item Encuentre el número que falta en cada caso:
        \begin{enumerate}[label=\alph*)]
          \item $\_\_\_ \times (-7) = 56$
          \item $(-48) \div \_\_\_ = -6$
        \end{enumerate}

  \item Primero decida el signo del resultado (contando los factores negativos) y luego calcule:
        \[
          (-5) \times (-4) \times (-3) \div (-2) \div (-5)
        \]
\end{enumerate}

\newpage

% ============================================================
\section{Operaciones Combinadas}
% ============================================================

\begin{enumerate}[leftmargin=*]
  \item Resuelva:
        \[
          -5 \times [4 - (-3) \times 2] \div (-2)
        \]

  \item Resuelva, evaluando de adentro hacia afuera:
        \[
          [-12 \div (-3) - (-4) \times 2] \div (-2) - (-1)
        \]

  \item Resuelva:
        \[
          \{[6 - (-2)] \times (-3) + (-12)\} \div (-6)
        \]

  \item Resuelva la operación multinivel:
        \[
          -2 \times \{-5 + 3 \times [-4 - (-12 \div 3 + 1)]\} - (-10)
        \]
\end{enumerate}

\newpage

% ============================================================
\section{Problemas con Números Enteros}
% ============================================================

\begin{enumerate}[leftmargin=*]
  \item En la sierra, a las 6 de la mañana la temperatura era de $-8\,^{\circ}$C. A lo largo del día subió $15\,^{\circ}$C, luego bajó $22\,^{\circ}$C, subió $7\,^{\circ}$C y finalmente bajó $3\,^{\circ}$C. \textquestiondown Cuál fue la temperatura final?

  \item Un ascensor parte del piso $-2$ (sótano). Sube $9$ pisos, baja $5$, sube $12$, baja $7$ y finalmente baja $6$. \textquestiondown En qué piso termina?

  \item Martín debía $24$ Bs. a su amigo. El lunes le pagó $9$ Bs., el martes $5$ Bs. y el miércoles $12$ Bs. Representa la deuda como un número entero y determina cuánto debe al final.

  \item Un buzo en una competencia parte de la superficie ($0$ m). Desciende $4$ metros por minuto durante $3$ minutos, luego asciende $9$ metros y después desciende $2$ metros por minuto durante $5$ minutos. \textquestiondown A qué profundidad termina?
\end{enumerate}

\newpage

% ============================================================
\ifmwclaves
  \section{Clave de Respuestas}
  % ============================================================

  \subsection*{Suma y Resta en $\Z$ (Ejercicios 1--4)}

  \begin{enumerate}[leftmargin=*, label=\textbf{\arabic*.}]
    \item $|-7| - |-11| + |7| = 7 - 11 + 7 = 3$
    \item $(3 - 9) = -6 \Rightarrow -(-6) = 6$; \ $(5 - 12) = -7 \Rightarrow -(-7) + 4 = 11$; \ $6 - 11 = -5$
    \item $15 - 8 - 6 - 4 = -3$
    \item $-10 + x + 5 = -18 \Rightarrow x = -13$
  \end{enumerate}

  \subsection*{Multiplicación y División en $\Z$ (Ejercicios 1--4)}

  \begin{enumerate}[leftmargin=*, label=\textbf{\arabic*.}]
    \item $(-4) \times (-6) = 24$; $24 \div (-3) = -8$; $-8 \times 2 = -16$
    \item $(-9) \times 8 = -72$; $-72 \div (-6) = 12$; $12 \div (-4) = -3$
    \item a) $-8$ \quad b) $8$
    \item Cinco factores negativos (cantidad impar) $\Rightarrow$ resultado negativo.
          $(-5) \times (-4) = 20$; $20 \times (-3) = -60$; $-60 \div (-2) = 30$; $30 \div (-5) = -6$
  \end{enumerate}

  \subsection*{Operaciones Combinadas (Ejercicios 1--4)}

  \begin{enumerate}[leftmargin=*, label=\textbf{\arabic*.}]
    \item $(-3) \times 2 = -6$; $4 - (-6) = 10$; $-5 \times 10 = -50$; $-50 \div (-2) = 25$
    \item $-12 \div (-3) = 4$; $(-4) \times 2 = -8$; $4 - (-8) = 12$; $12 \div (-2) = -6$; $-6 - (-1) = -5$
    \item $6 - (-2) = 8$; $8 \times (-3) = -24$; $-24 + (-12) = -36$; $-36 \div (-6) = 6$
    \item $(-12 \div 3 + 1) = -3$; $[-4 - (-3)] = -1$; $3 \times (-1) = -3$; $-5 + (-3) = -8$;
          $-2 \times (-8) = 16$; $16 - (-10) = 26$
  \end{enumerate}

  \subsection*{Problemas con Números Enteros (Ejercicios 1--4)}

  \begin{enumerate}[leftmargin=*, label=\textbf{\arabic*.}]
    \item $-8 + 15 = 7$; $7 - 22 = -15$; $-15 + 7 = -8$; $-8 - 3 = -11\,^{\circ}$C
    \item $-2 + 9 = 7$; $7 - 5 = 2$; $2 + 12 = 14$; $14 - 7 = 7$; $7 - 6 = 1$ \ (piso $1$)
    \item $-24 + 9 + 5 + 12 = 2$ \ $\Rightarrow$ la deuda queda en $0$ y le sobran $2$ Bs.
    \item Descenso inicial: $3 \times (-4) = -12$; $-12 + 9 = -3$;
          descenso final: $5 \times (-2) = -10$; $-3 - 10 = -13$ m
  \end{enumerate}

  \vspace{10pt}
  \begin{cuadrosolucion}
    \textbf{\textquestiondown Cómo te fue?} Si manejaste bien la regla de los signos y la jerarquía de las operaciones, \textquestiondown listo para el Tema 3! Concentra tu atención en los \emph{signos dobles} y en el \emph{orden} de las operaciones combinadas.
  \end{cuadrosolucion}

\fi
\end{document}
